% chapter structure

% ********************
\section{Introduction}
% ********************

Un rapport est par définition un document structuré contenant des informations à faire passer.
 La structure, indépendamment du contenu permet au lecteur 
de mieux suivre le document et d'aller rapidement  chercher les informations intéressantes pour lui.

La structure d'un rapport est le squelette qui décrit les éléments dans lesquels seront ajoutés
plus tard le contenu (texte, tableau, figures, images).

Les éléments sont des pages particulières, des chapitres ou des sections, voir des éléments inférieures comme des sous-sections et des paragraphes. 
La bibliographie par exemple est un élément de structure.

La structure générale sera décrite brièvement dans un premier temps.
Ensuite des éléments de contenu ou de format seront fournis pour chaque structure.
Enfin on conclura sur la numérotation d'un document.

% *****************************
\section{Un rapport structuré}
% *****************************

Une revue complète des éléments de structures se trouvent dans cette section.
La structure la plus importante est ici le chapitre, mais comme cela a été indiqué dans le chapitre 'introduction',
 si le document est trop court (inférieur à 20 pages), il vaut mieux que la structure principale reste la section.
 Dans les revues internationales on trouve  des articles de plus de 50 pages constitués de section uniquement, mais cela est imposé par l'éditeur.  

Les éléments sont classés par ordre probable d'apparition dans le document:

\begin{enumerate}
\item une page de \textbf{titre}
\item un \textbf{résumé} obligatoire d'une page seulement.
\item un \textbf{sommaire}
\item une \textbf{liste des figures} avec la légende (optionnelle)
\item une \textbf{liste des tableaux} avec la légende (optionnelle)
\item des \textbf{remerciements} (optionnels)
\item une \textbf{nomenclature} (optionnelle suivant le contenu)
\item un chapitre \textbf{Introduction} décrivant les objectifs du projet et indiquant brièvement les
méthodologies, les tâches effectuées et les principaux résultats attendus.
\item un chapitre ou une section sur les \textbf{partenaires} et l'\textbf{organisation générale du projet}\footnote{La gestion de projet} incluant par exemple une \textbf{liste de livrables} (documents à fournir comme preuve du travail).
\item  Un chapitre ou une section \textbf{bibliographie} décrivant l'état de l'art sur la thématique du projet. Elle peut être inclus dans l'introduction ou dans une section spéciale.
\item des chapitres structurés en sections décrivant le travail effectué, les méthodologies, les outils, les résultats et des analyses scientifiques ou techniques.
\item Un chapitre ou une annexe ou une section sur les \textbf{aspects financiers} du projet
\item Un chapitre indiquant \textbf{conclusions et perspectives}
\item Un ou plusieurs  chapitres \textbf{annexes} 
\end{enumerate}

Certains éléments ne sont pas toujours pertinents, cela dépend du contexte. L'ordre est aussi modifiable tout en respectant la logique.

Il ne faut pas hésiter à regarder des exemples dans des livres ou des rapports officiels, des thèses  ou 
des rapports scientifique de la NASA par exemple.


La figure \ref{fig:structure_rapport} résume les recommandations indiquées précédemment.

\begin{figure}[htb]
\centering
\includegraphics[height=10cm]{Report_Structure.pdf}
\caption{Structure d'un rapport.}
\label{fig:structure_rapport}
\end{figure}

% *******************************
\section{Numérotation des pages}
% *******************************


Voici quelques indications de mise en forme  qu'il faut essayer de suivre:

\begin{itemize}

\item un chapitre commence toujours sur une page de droite, qui porte un numéro impair. Ne pas hésiter à laisser une page blanche avant si ce n'est pas le cas.

\item  démarrer la page 1 sur le chapitre Introduction.

\item Les pages avant l'introduction ont des numéros romains ou  ne sont pas numérotées

\item Les pages des annexes sont numérotées normalement.

\end{itemize}


%***********************
\section{Page de titre}
%***********************
On rappelle ici les informations  qu'on trouve généralement sur une page de titre :
\begin{itemize}
\item des logos de l'université et de l'entreprise ou du laboratoire
\item le nom du diplôme
\item les mots "rapport de stage", "rapport de projet", ...
\item un titre avec éventuellement un sous-titre
\item un ou plusieurs noms des auteurs du rapport, avec leur affiliation si tout le monde ne vient pas du
même organisme ou entreprise
\item Le nom de l'encadrant (souvent)
\item une date
\item une ou des images d'illustrations, en rapport avec le thème du rapport
\item éventuellement on peut y trouver des adresses.
\end{itemize}



% *************************
\section{Sommaire }
% *************************

\textit{(contents en anglais)}

Un sommaire ne doit pas dépasser que quelques pages. 

Il faut définir le niveau des sections qui seront présentes avec en quelque sorte un ordre:

\begin{enumerate}
\item chapitre
\item section
\item sous section
\item sous sous section
\item paragraphe
\end{enumerate}

Le paragraphe est situé au niveau 5, le chapitre au niveau 1.

Ainsi on peut décider que le sommaire s'arrêtera au niveau 2 ou niveau 3  ou 4 par exemple.


\Important{
Dans un document il ne faut pas de sections ou sous sections dites mortes:
\begin{itemize}
\item Un rapport ne peut pas être constitué d'un unique chapitre
\item Un chapitre ne peut pas être constitué d'une unique section
\item Une section ne peut pas être constituée d'une unique sous-section
\item etc ...
\end{itemize}

Dans ce cas il faut supprimer l'élément unique, car cela ne sert à rien de l'avoir introduit.

En regardant la table des matières on trouve tout de suite les éléments uniques donc inutiles.
}

Pour qu'un rapport soit harmonieux il faut une structure assez similaire d'un chapitre à l'autre, avec approximativement le même nombre 
de sections. Cela se voit très rapidement dans le sommaire.


% *********************
\section{Nomenclature}
% *********************
\label{sec:nomenclature}

La nomenclature permet de rappeler les notations utilisées dans le rapport.
\rouge{Cette section est optionnelle,  cela devient nécessaire afin de ne pas perdre le lecteur 
lorsqu'un grand nombre de notations sont employées.}
 

On peut analyser des exemples dans des rapports techniques ou scientifiques qu'on trouve sur internet (par exemple les papiers des conférences AIAA)

\vs{0.5}

Elle est en général décrite par un tableau contenant divers champs ou sections qui décrivent: 
\begin{itemize}
\item les variables en lettres romanes
\item les variables en lettres grecques
\item les quantités physiques ou sans dimensions
\item les acronymes
\item la notation pour les indices
\item la notation pour les exposants
\item la notation avec des symboles surlignés et soulignés comme par exemple
$\tilde x$, $\dot x$ ou $\underline{y}$.
\end{itemize}

La nomenclature n'a pas toujours besoin d'être exhaustive. 

Ainsi si un symbôle n'intervient qu'une fois dans le document on peut se contenter de 
l'expliquer dans le texte.


Il faut adapter la présentation en fonction du nombres d'éléments à présenter.

Un exemple de nomenclature se trouve au début de ce rapport.


% *********************************
\section{Introduction - Conclusions}
% **********************************

Ces deux parties sont toujours à écrire en dernier, dans l'ordre, conclusion, puis introduction.

Dans cette section on se réfère aux chapitres introductions et conclusions, mais il faut aussi comprendre que chaque chapitre
doit avoir une section introduction et conclusion.

%*************************
\subsection{Introduction}
%*************************

On pourra se reporter à la section \ref{sec:introduction} de ce document.

Typiquement, cette section est un exemple typique de section d'un document qui ne sert à rien et qui doit être supprimée.
Ou bien il faut prendre des éléments du chapitre \ref{sec:introduction} pour les amener ici.

Il n'en demeure pas moins que l'introduction est nécessaire et obligatoire, C'est la première partie réellement lues du document.
Elle est souvent à l'image du reste du document, il faut donc qu'elle soit de qualité.

%*************************
\subsection{Conclusion}
%*************************

A priori c'est la dernière partie qui est lue avec intérêt. Elle donnera l'impression finale sur le document. Elle doit aussi être de qualité.
Naturellement elle est obligatoire.

Certaines personnes, comme moi, lisent la conclusion avant l'introduction, car c'est la seule partie où on peut, où on doit apprendre
des choses en quelques lignes, puisque la conclusion rappelle la problématique, la méthodologie, indique les principaux résultats et dressent
des conclusions et des perspectives.


\rouge{Typiquement comme il n'y avait pas grand chose à ajouter à cette section, on a regroupé deux sections sous une seule en créant deux sous-sections.}


%**************************
\section{Bibliographie}
%**************************

C'est habituellement l'élément où la qualité est la plus mauvaise, si on trouve la section dans un rapport.

Pourtant la bibliographie indique clairement comment a été effectué le travail, sur quelles informations ou données il est basé. C'est aussi une
partie souvent regardée avant de lire le document. 

La bibliographie suit un format rigoureux avec des informations dans un ordre défini qui doit être respecté pour toutes les références.

Toute référence doit être complète afin d'être trouvée facilement. Les informations minimales sont le nom des auteurs avec les initiales des prénoms,
le titre de la référence et l'année de publication.
Ensuite on doit trouver :
\begin{itemize}
\item si c'est un article :  le nom de la revue, le volume et le numéro, le numéro des pages est optionnel, le DOI devient de plus en plus obligatoire.
\item si c'est un rapport de thèse ou de master : le type de rapport, l'université, éventuellement un numéro de référence
\item si c'est un rapport scientifique ou technique : le type du rapport, le nom de l'organisme, un numéro de référence
\item si c'est un livre : le nom de l'éditeur, le DOI devient de plus en plus obligatoire.
\item si c'est un papier d'une conférence: le numéro du papier, le nom complet de la conférence, le lieu et la date.  
\end{itemize}

On trouve de plus en plus souvent une webographie, c'est à dire une bibliographie liée à des sites webs. Il faut alors donner l'adresse complète du site qui permet de trouver facilement l'élément cité (une image, un film, un texte, ...).

On peut mettre la webographie dans la bibliographie.


\Attention{Quand on ne trouve rien à mettre dans la bibliographie, il faut se poser des questions sur la façon dont on a travaillé. 
Quelque soit le travail demandé on se base toujours sur des connaissances qu'on a trouvées quelque part.

Si une entreprise ne souhaite pas qu'on ajoute les rapports internes pour raisons de confidentialité, alors il faut soit mettre uniquement le titre avec rapport interne, et année, sans plus de précision, soit justifier dans le texte le choix de ne pas mettre de bibliographie.

\textbf{Il est extrêmement rare que l'absence de bibliographie soit justifiable, même dans un contexte industriel.}
}


% ***********************
\section{Les annexes}
% ***********************


Les annexes permettent d'ajouter des informations complémentaires au rapport qui ne sont pas absolument nécessaires, cela puet être :
\begin{itemize}
\item des démonstrations mathématiques
\item des détails sur tels ou tels méthodologies
\item des figures ou des images
\item des tableaux de chiffres ou autres bases de données utiles ou utilisées
\item le développement d'un point particulier jugé suffisamment important pour être mentionné mais pas suffisamment important pour appartenir au corps du texte.
\end{itemize}


Attention aux points suivants:
\begin{enumerate}
	\item On ne trouve jamais plus de 4 ou 5 annexes, et bien qu'il faut considérer qu'ils ne seront pas nécessairement lus ou regardés, 
	il faut rédiger à minima  quelques phrases bien écrites. 

	\item Il faut éviter des photocopies ou documents illisibles car trop petits ou de très mauvaises qualités. 

 	\item L'annexe de doit pas être trop épais (moins de 10 pages habituellement)
 	
 	\item Il ne faut pas imprimer des codes informatiques en annexe, sauf si c'est exigé. Dans ce cas, il faut utiliser un traitement de texte pour minimiser la place,
 	par exemple en utilisant plusieurs colonnes, en diminuant la taille des caractères à 10 ou 9 points. Il faut mieux privilégier l'algorithme du code 
 	       et donner les codes sur une  clé USB avec le rapport ou fournir un lien de téléchargement.
\end{enumerate}


% **************************
\section{Les remerciements}
% **************************

Cette page bien que très personnelle ne doit pas vous porter préjudice si quelqu'un qui ne vous connait pas la lit.

Il faut qu'elle soit écrite dans un beau français ou bel anglais, et se contenter de remerciements classiques, des personnes qui vous ont aidés
directement ou indirectement, sans mentions de faits
ou d'actions qui n'ont rien à voir avec le projet, ou bien cela doit être écrit avec beaucoup de sobriété et de délicatesse.
D'une certaine façon on devine à la lecture des remerciements comment le stagiaire, pour le cas d'un rapport de stage, c'est intégré à l'entreprise.

Il faut éviter de mêler la religion à la science sauf si vous voulez ouvrir un nouveau débat. Mais le retour pourra être catastrophique si vous
devez envoyer votre rapport à un futur employeur.

Dans beaucoup de rapports scientifiques, on remercie les financeurs du projet ou les centres de calculs qui ont permis de mener le travail à bien.
Cela est à discuter avec l'encadrement.
C'est un point prend de plus en plus d'importance à l'ère du numérique ou tout est référencé.

En conclusion, il ne faut pas négliger cette partie qui peut au final participer à améliorer ou pas l'évaluation du travail et du rapport, juste par quelques mots
bien choisis ou mal à propos.